% Official CC BY4.0 published HTML A2; original math alttext and author text, literal original numbering.


\par\medskip\noindent\textbf{Appendix B PBW, proof that monomials span the space}\par




In this section, we give an elementary proof of the implication in Theorem 4.9 that we use. We include this standard result in this appendix, because it is the main result that enables our algorithmic speedup compared to a naive implementation.





Recall that the universal enveloping algebra $\mathcal{U}(\mathfrak{g})$ is the quotient of the tensor algebra $\mathcal{T}(\mathfrak{g})$ by the ideal $\mathcal{I}(\mathfrak{g})=\{X\otimes Y-Y\otimes X-[X,Y]:X,Y\in\mathfrak{g}\}$. The quotient induces a filtration $\mathcal{U}(\mathfrak{g})_{\leq m}\subseteq\mathcal{U}(\mathfrak{g})_{\leq(m+1)}$ and the elements of $\mathcal{U}(\mathfrak{g})_{\leq m}$ are represented by noncommutative polynomials of degree at most $m$ in the elements of $\mathfrak{g}$.





\par\medskip\noindent\textbf{Theorem B.1 (One direction of the PBW Theorem).}\par



Let $X_{1},\ldots,X_{r}$ be a basis for a Lie algebra $\mathfrak{g}$. Then



\[\begin{aligned}
\displaystyle\mathcal{B}_{m}\coloneqq\big{\{}X_{i_{1}}X_{i_{2}}\cdots X_{i_{\ell}}\mid i_{1}\leq i_{2}\leq\cdots\leq i_{\ell}\text{ and }\ell\leq m\}
\end{aligned}\]


spans $\mathcal{U}(\mathfrak{g})_{\leq m}$.






\par\medskip\noindent\textbf{Proof.}\par




It is clear that $\mathcal{T}(\mathfrak{g})_{m}$ is spanned by all elements $X_{i_{1}}\otimes\cdots\otimes X_{i_{m}}$; hence it suffices to prove that any such element is equivalent, modulo $\mathcal{I}(\mathfrak{g})$, to a linear combination of elements of $\mathcal{B}_{m}$. We prove this statement by induction on $m$.





Let $X_{j_{1}}\otimes\cdots\otimes X_{j_{m}}\in\mathcal{T}(\mathfrak{g})_{m}$. For every $a=1,\ldots,m-1$, we have $X_{j_{a}}\otimes X_{j_{a+1}}-X_{j_{a+1}}\otimes X_{j_{a}}-[X_{j_{a}},X_{j_{a+1}}]\in\mathcal{I}(\mathfrak{g})$; in particular $X_{j_{a}}\otimes X_{j_{a+1}}=X_{j_{a+1}}\otimes X_{j_{a}}+Y$ where $Y$ is a linear combination of elements $\mathcal{I}(\mathfrak{g})$ and of $\mathcal{T}(\mathfrak{g})_{\leq 1}$. We obtain



\[\begin{aligned}
\begin{split}X_{j_{1}}\otimes\cdots\otimes X_{j_{\ell}}=&X_{j_{1}}\otimes\cdots\otimes X_{j_{a-1}}\otimes X_{j_{a+1}}\otimes X_{j_{a}}\otimes X_{j_{a+2}}\otimes\cdots\otimes X_{j_{\ell}}+Y\end{split}\qquad\text{(18)}
\end{aligned}\]


where $Y$ is a linear combination of elements of $\mathcal{I}(\mathfrak{g})$ and of $\mathcal{T}(\mathfrak{g})_{\leq m-1}$.





Let $\pi\in\mathfrak{S}_{m}$ be a permutation such that $(j_{\pi(1)}\cdots j_{\pi(m)})$ is nondecreasing. Write $\pi$ as a composition of elementary transpositions $\pi=\sigma_{s}\cdots\sigma_{1}$, where $\sigma_{q}=(a,a+1)$ for some $a=1,\ldots,m-1$. Applying Equation 18 $s$ times we obtain



\[\begin{aligned}
X_{j_{1}}\otimes\cdots\otimes X_{j_{m}}=X_{j_{\pi(1)}}\otimes\cdots\otimes X_{j_{\pi(m)}}+U
\end{aligned}\]


where $U\in\mathcal{I}(\mathfrak{g})+\mathcal{U}_{\leq(m-1)}$. Passing to the quotient modulo $\mathcal{I}(\mathfrak{g})$, we obtain the following identity in $\mathcal{U}(\mathfrak{g})_{\leq m}$:



\[\begin{aligned}
X_{j_{1}}\cdots X_{j_{m}}=X_{j_{\pi(1)}}\cdots X_{j_{\pi(m)}}+U
\end{aligned}\]


with $U\in\mathcal{U}(\mathfrak{g})_{\leq(m-1)}$. Apply the induction hypothesis to $U$ to conclude.
$\hfill\blacktriangleleft$




